\chapter{Anexo I: Código de la Solución}
\label{cap:anexo_codigo_solucion}

\section{Código de WinCC}
\definecolor{codegreen}{rgb}{0,0.6,0}
\definecolor{codegray}{rgb}{0.5,0.5,0.5}
\definecolor{codepurple}{rgb}{0,0,0.82}
\definecolor{backcolour}{rgb}{1,1,1}

\lstset{
    language=C,
    backgroundcolor=\color{backcolour},   
    commentstyle=\color{codegreen},
    keywordstyle=\color{magenta},
    numberstyle=\tiny\color{codegray},
    stringstyle=\color{codepurple},
    basicstyle=\ttfamily\scriptsize,
    breakatwhitespace=false,         
    breaklines=true,                 
    captionpos=b,                    
    keepspaces=true,                 
    numbers=left,                    
    numbersep=10pt,                  % Reducido a 10pt para que no se desplace demasiado al quitar el frame
    frame=single,                % Añade un borde muy fino
    framerule=0.2pt,
    xleftmargin=5pt,                % Añade margen izquierdo al bloque de color de fondo
    xrightmargin=5pt,                % Añade margen derecho al bloque de color de fondo
    aboveskip=0.75\baselineskip,
    belowskip=0.25\baselineskip,
    showspaces=false,                
    showstringspaces=false,
    showtabs=false,                  
    tabsize=2
}

\section{Código fuente}
\subsection{Filtrado de turbinas por estado y subestación}
\begin{lstlisting}
void AplicarFiltroPosicion(){

    int selValue, parkValue;
    int CANT_MOLINOS, i;
    int tagRun, tagEmg, tagPark, state;
    int match;

    // Variables de cuadricula
    int startX = 10;
    int startY = 0;
    int offsetX = 475;
    int offsetY = 220;
    int maxCols = 4;

    int currentX;
    int currentY;
    int colCount = 0;
    //Nombres de tags
    char tagNameRun[100];
    char tagNameEmg[100];
    char tagNamePark[100];
    char FP[50];

    // Guarda IDs de molinos que cumplen el filtro
    int validList[120];
    int validCount = 0;

    char* prefijoActual;

    currentX = startX;
    currentY = startY;
    // Obtener valores de filtros
    selValue = GetTagWord("HMI_Filtro_Molinos");
    parkValue = GetTagWord("HMI_Filtro_Parques");
    prefijoActual = "WTG";
    CANT_MOLINOS = 120;
    // Evaluar cada molino
    for (i = 1; i <= CANT_MOLINOS; i++) {
        sprintf(tagNameEmg, "%s%d.WTUR.TurSt.stVal", prefijoActual, i);
        sprintf(tagNamePark, "%s%d.LLNO.InSSet.setSrc", prefijoActual, i);

        tagEmg = GetTagWord(tagNameEmg);
        tagPark = GetTagWord(tagNamePark);

        switch (tagEmg) {
            case 0:
                state = 2;
                break;
            case 1:
            case 4:
                state = 3;
                break;
            case 2:
                state = 4;
                break;
            case 3:
                state = 5;
                break;
            case 5:
            case 7:
                state = 6;
                break;
            case 6:
                state = 7;
                break;
            default:
                if (tagEmg >= 10 && tagEmg <= 19) {
                    state = 8;
                } else {
                    state = -1;
                }
                break;
        }

        match = 0;
        if (parkValue == 1 || parkValue == tagPark) {
            if (selValue == 1 || selValue == state) {
                match = 1;
            }
        }

        if (match == 1) {
            validList[validCount] = i;
            validCount++;
        }
    }
    // Visualizar molinos que cumplen el filtro
    for (i = 1; i <= CANT_MOLINOS; i++) {
        sprintf(FP, "PW_%d", i);
        SetPropBOOL("20_WTG_Turbine_Overview.PDL", FP, "Visible", FALSE);
    }

    for (i = 0; i < validCount; i++) {
        int molinoID = validList[i];
        sprintf(FP, "PW_%d", molinoID);

        SetPropBOOL("20_WTG_Turbine_Overview.PDL", FP, "Visible", TRUE);
        SetPropDouble("20_WTG_Turbine_Overview.PDL", FP, "Left", currentX);
        SetPropDouble("20_WTG_Turbine_Overview.PDL", FP, "Top", currentY);

        colCount++;
        if (colCount >= maxCols) {
            currentX = startX;
            currentY += offsetY;
            colCount = 0;
        } else {
            currentX += offsetX;
        }
    }
}
\end{lstlisting}
\subsection{Filtrado de parques por estado y subestación}
\begin{lstlisting}
#include "apdefap.h" 
void AplicarFiltroParque(){ 

    //Variables
   
    int selValue, parkValue; 
    int CANT_VENTANAS = 22; 
    int i, k; 
    int match; 

    int tagParkBase;    // Parque del molino principal (TagPrefix)
    int currentPark;    // Parque del molino que estamos evaluando en el bucle k
    int baseID;         // ID numerico del molino principal 
    int currentID;      
    int rawStatus; 
    int state; 

    int startX = 10;    
    int startY = 0;     
    int offsetX = 475;  
    int offsetY = 225;  
    int maxCols = 4; 
    
    int currentX = startX; 
    int currentY = startY; 
    int colCount = 0;  

    char FP[50];  
    char* prefijoVentana; 
    char tagName[150]; 

    //Leer valores de filtros
    selValue = GetTagWord("HMI_Filtro_Molinos"); 
    parkValue = GetTagWord("HMI_Filtro_Parques"); 

  
    // Evaluar cada ventana de imagen (cada una con un parque base diferente)

    for (i = 1; i <= CANT_VENTANAS; i++) { 
        match = 0; 
        
        sprintf(FP, "PK_%d", i); 
        
        // LEEMOS EL TAG PREFIX
        prefijoVentana = GetPropChar("21_WTG_Park_Overview.PDL", FP, "TagPrefix"); 
        
        if (prefijoVentana != NULL && strlen(prefijoVentana) > 0) { 
            
            // Leemos el parque del molino base para referencia
            sprintf(tagName, "%s.LLNO.SubSt.setSrc", prefijoVentana); 
            tagParkBase = GetTagWord(tagName); 
            
            // Leemos el ID numerico base
            sprintf(tagName, "%s.LLNO.UnitID.setSrc", prefijoVentana); 
            baseID = GetTagWord(tagName); 

            if (baseID > 0) {
                // BUCLE K: Comprobamos el molino base (0), el +1 y el +2
                for (k = 0; k < 3; k++) { 
                    currentID = baseID + k; 

                    // 1. Obtener parque del molino actual (k) para ver si esta en el parque filtrado
                    if (k == 0) {
                        currentPark = tagParkBase;
                    } else {
                        sprintf(tagName, "WTG%d.LLNO.SubSt.setSrc", currentID); 
                        currentPark = GetTagWord(tagName);
                    }

                    // FILTRO DE PARQUE: El molino 'k' debe pertenecer al parque seleccionado
                    if (parkValue == 1 || parkValue == currentPark) {
                        
                        // 2. Leer estado del molino 'k'
                        if (k == 0) {
                            sprintf(tagName, "%s.WTUR.TurSt.stVal", prefijoVentana);
                        } else {
                            sprintf(tagName, "WTG%d.WTUR.TurSt.stVal", currentID);
                        }
                        rawStatus = GetTagWord(tagName); 

                        // 3. Maquina de estados interna 
                        switch (rawStatus) { 
                            case 0: state = 2; break; // Stopped 
                            case 1: 
                            case 4: state = 3; break; // Starting o Shutting down 
                            case 2: state = 4; break; // Running 
                            case 3: state = 5; break; // Paused 
                            case 5: 
                            case 7: state = 6; break; // Emergency stop o Fault 
                            case 6: state = 7; break; // Maintenance 
                            default: 
                                if (rawStatus >= 10 && rawStatus <= 19) { 
                                    state = 8; // Communication Error 
                                } else { 
                                    state = -1; 
                                } 
                                break; 
                        } 

                        //filtro estado
                        if (selValue == 1 || selValue == state) { 
                            match = 1;  
                            break; // Si uno de los tres cumple, ya no hace falta mirar mas
                        } 
                    }
                } 
            } 
        } 
        
       
        // Visualizacion de ventanas: Si 'match' es 1, esta ventana se muestra en la cuadricula, si no se oculta
    
        if (match == 1) { 
            SetPropDouble("21_WTG_Park_Overview.PDL", FP, "Left", currentX); 
            SetPropDouble("21_WTG_Park_Overview.PDL", FP, "Top", currentY);  
            SetPropBOOL("21_WTG_Park_Overview.PDL", FP, "Visible", TRUE); 
            
            colCount++; 
            if (colCount >= maxCols) { 
                currentX = startX; 
                currentY += offsetY; 
                colCount = 0; 
            } else { 
                currentX += offsetX; 
            } 
        } else { 
            SetPropBOOL("21_WTG_Park_Overview.PDL", FP, "Visible", FALSE); 
        } 
    } 
}
\end{lstlisting}
\subsection{Navegación menu}
\begin{lstlisting}
Sub MenuNavigation(Byval menuItem)
Dim myScreen : myScreen =menuItem.userData
HMIRuntime.Screens(BaseScreenName).ScreenItems("img_GUI").TagPrefix=""
HMIRuntime.Screens(BaseScreenName).ScreenItems("img_GUI").PictureName=myScreen

If myScreen ="13_SYS_Alarms" Then 
	HMIRuntime.Screens(BaseScreenName).ScreenItems("img_alarm").PictureName="DiagLog"
Else
	HMIRuntime.Screens(BaseScreenName).ScreenItems("img_alarm").PictureName="AlarmLog"
End If


End Sub
\end{lstlisting}
\section{Código de IIoT}

\begin{lstlisting}[language=bash]
sudo apt update
sudo apt install docker.io docker-compose -y
sudo systemctl start docker
sudo systemctl enable docker
\end{lstlisting}

% Tip: He cambiado temporalmente a lenguaje=YAML para que el resaltado del compose sea correcto
\begin{lstlisting}[language=XML]
version: '3.8'

services:
  mosquitto:
    image: eclipse-mosquitto:latest
    container_name: broker_mqtt
    ports:
      - "1883:1883"
    volumes:
      - ./mosquitto/config/mosquitto.conf:/mosquitto/config/mosquitto.conf
    restart: always

  influxdb:
    image: influxdb:2.7
    container_name: base_datos
    ports:
      - "8086:8086"
    volumes:
      - ./influxdb_data:/var/lib/influxdb2
    restart: always

  nodered:
    image: nodered/node-red:latest
    container_name: puente_nodered
    ports:
      - "1880:1880"
    volumes:
      - ./nodered_data:/data
    restart: always

  grafana:
    image: grafana/grafana:latest
    container_name: graficos_web
    ports:
      - "3000:3000"
    volumes:
      - ./grafana_data:/var/lib/grafana
    restart: always

\end{lstlisting}

\begin{lstlisting}[language=bash]
sudo docker-compose up -d
\end{lstlisting}

\vspace{1cm}


